\section{Summary of Key Findings}
\label{sec:summary}

The literature has moved from single-environment causal discovery to methods that
use known environments, distribution shifts, or joint causal clustering. This
section summarises the main trends and limitations and states the gaps addressed by
this thesis.

\subsection{Methodological Trends}

Constraint-based methods use conditional independence, while score-based methods
optimize a graph score such as BIC. Functional causal models, including LiNGAM and
VAR-LiNGAM, use stronger assumptions about the noise and functional form
\cite{assaad2022survey,shimizu2006lingam,hyvarinen2010estimation}.

In heterogeneous settings, invariance and distribution shifts are the main guiding
ideas. They appear in ICP, CD-NOD, CD-NOTS, IRM, and JCI
\cite{peters2016causal,huang2020causal,sadeghi2025causal,arjovsky2019irm,mooij2020jci}.
CASTOR and DyCAST model temporal changes as piecewise-constant or continuously
evolving structures \cite{rahmani2025castor,cheng2025dycast}. Causal clustering
methods such as SSCM, CCSL, HCL, and MCD jointly estimate assignments and causal
models \cite{huang2019specific,chen2023ccsl,li2025hcl,varambally2024mcd}.

\subsection{Consistent Findings}

The reviewed studies support five conclusions:
\begin{enumerate}[nosep]
\item Causal structure can differ substantially across regimes. A pooled model can
produce a graph that represents none of them
\cite{huang2020causal,chen2023ccsl,varambally2024mcd,li2025hcl,rahmani2025castor}.
\item Joint estimation of assignments and causal structures is generally better than
a sequential pipeline \cite{huang2019specific,chen2023ccsl,li2025hcl,rahmani2025castor}.
\item Causal-structure similarity can be a better clustering signal than feature
distance or correlation \cite{chen2023ccsl}.
\item Real data often violates stationarity, as shown by CD-NOTS and CASTOR
\cite{sadeghi2025causal,rahmani2025castor}.
\item No method is universally best. Performance depends on whether its assumptions
match the data-generating process \cite{assaad2022survey}.
\end{enumerate}

\subsection{Contradictions and Limitations}

ICP and IRM require known environment partitions. CD-NOD and CD-NOTS require an
auxiliary domain or time index. CCSL assumes stationarity within clusters, while
CASTOR handles sequential changes but assumes piecewise-constant regimes and
Gaussian equal-variance noise. These requirements limit their use when the hidden
state and the number of regimes are both unknown.

There is also a trade-off between flexibility and identifiability. CCSL has formal
guarantees under linear non-Gaussian noise, but that assumption may be restrictive.
HCL and MCD use more flexible models but provide weaker or less developed formal
guarantees. Evaluation is also inconsistent: some studies report clustering quality
without graph quality, while others do the reverse. A common benchmark is still
needed. Internal incoherency scores are promising for detecting assumption
violations, but they have not been adapted to causal clustering
\cite{faltenbacher2025internal}.

\subsection{Research Gaps}
\label{ssec:gaps}

The main gaps relevant to this thesis are:
\begin{itemize}[leftmargin=1.4em]
\item \textbf{No formal invariance test for discovered boundaries.} CASTOR and
DyCAST use global optimization scores, but do not attach a direct hypothesis test
to each proposed boundary. LIRA addresses this with a Chow test followed by an
independent BIC decision.
\item \textbf{No formal test of structural difference.} Existing causal-clustering
methods generally do not test whether two learned regimes differ beyond finite-sample
variation.
\item \textbf{Limited scalability evidence.} Most causal-clustering studies use
small synthetic graphs. Comparable large-graph evaluations remain limited.
\item \textbf{No standard causal-clustering benchmark.} The field lacks agreed
datasets and evaluation measures for comparing regime recovery and graph recovery.
\item \textbf{Limited observational identifiability theory.} Formal guarantees for
mixtures of SCMs from observational i.i.d. data remain limited beyond the linear
non-Gaussian case of CCSL.
\item \textbf{Within-regime nonstationarity.} Most methods assume that each regime
is stationary, even though the mechanism may continue to change inside a regime.
\item \textbf{Assumption checks without ground truth.} Tools for detecting violated
causal assumptions have not been developed fully for causal clustering.
\end{itemize}

This thesis addresses the first two gaps directly and evaluates regime recovery,
graph recovery, and scalability on controlled synthetic data. The remaining gaps,
including standard benchmarks, within-regime nonstationarity, and broader
assumption checks, are left for future work.
